\subsection{双线性型与半双线性型。秩。}

本节中，A 表示一个环（相应地一个配备了反自同构 J 的环），E 表示左 A-模，F 表示右 A-模（相应地左 A-模）。我们用左位似与右位似给 A 赋予 (A, A)-双模结构。这时，E × F 到双模 A 中的双线性（相应地关于 J 的右半双线性）映射，称为 E × F 上的双线性（相应地关于 J 的右半双线性）型。

当 E = F 时（这就意味着它是半双线性型），人们常把 E × F 上的一个半双线性型说成 E 上的一个半双线性型。

给定两个左 A-模 E 与 E′ 以及 E 与 E′ 上（关于 J）的两个半双线性型 \(\Phi\) 与 \(\Phi'\) ，称 \(\Phi\) 与 \(\Phi'\) 等价，如果存在 A-模 E 到 A-模 E′ 上的一个同构 \(u\) ，使得 \(\Phi'(u(x), u(y)) = \Phi(x, y)\) 对 E 中一切的 \(x\) 、\(y\) 成立；这时 \(\Phi\) 是 \(\Phi'\) 关于 \(u\) 与 \(u\) 的原像，而 \(\Phi'\) 是 \(\Phi\) 关于 \(u^{-1}\) 与 \(u^{-1}\) 的原像（第 2 小节）。

设 \(\Phi\) 为 \(E \times F\) 上的一个双线性型（其中 F 表示一个右 A-模）。线性映射 \(s_{\Phi}\) 与 \(d_{\Phi}\) （与 \(\Phi\) 相伴者，第 1 小节，定义 2）这时分别是 E 到 F 的对偶 \(F^{*}\) 中与 F 到 E 的对偶 \(E^{*}\) 中的映射。

于是，按定义，

\[
\Phi (x, y) = \left\langle x, d _ {\Phi} (y) \right\rangle = \left\langle y, s _ {\Phi} (x) \right\rangle .\tag{23}
\]

我们现在来定义与半双线性型相伴的线性映射。设 J 是 A 的一个反自同构，\(\Phi\) 是 E × F 上的一个（关于 J 的）（右）半双线性型（其中 F 表示一个左 A-模）；令 J′ = J \(^{-1}\) 。由下式定义的、F × E 到 A 中的映射 \(\Phi'\)

\[
\Phi^ {\prime} (y, x) = \Phi (x, y) ^ {J ^ {\prime}}
\]

\[
(x \in \mathrm{E}, y \in \mathrm{F})
\]

如容易看出的那样，是 F × E 上的一个（关于 J′ 的）（右）半双线性型。由第 2 小节（定义 5），半双线性型 Φ 与 Φ′ 分别等同于 E × F′ 上与 F × E′ 上的双线性型。与这两个双线性型相伴的映射 dΦ 与 dΦ′，分别称为半双线性型 \(\Phi\) 的右相伴与左相伴映射，记作 \(d_{\Phi}\) 与 \(s_{\Phi}\) 。于是，按定义我们有：

\[
\Phi (x, y) = \left\langle x, d _ {\Phi} (y) \right\rangle = \left\langle y, s _ {\Phi} (x) \right\rangle^ {J} \quad (x \in \mathrm{E}, y \in \mathrm{F}).\tag{24}
\]

于是，\(d_{\Phi}\) （相应地 \(s_{\Phi}\) ）是 \(\mathbf{F}^{\mathrm{j}}\) 到 \(\mathbf{E}^{*}\) 中（相应地 \(\mathbf{E}^{\mathrm{j}'}\) 到 \(\mathbf{F}^{*}\) 中）的线性映射，或者等价地说，是 \(\mathbf{F}\) 到 \(\mathbf{E}^{*}\) （相应地 \(\mathbf{E}\) 到 \(\mathbf{F}^{*}\) ）中的、关于 \(\mathbf{J}\) （相应地 \(\mathbf{J}'\) ）的半线性映射，如果我们把 \(\mathbf{J}\) （相应地 \(\mathbf{J}'\) ）看作环 \(\mathbf{A}^{0}\) （\(\mathbf{A}\) 的反环）到 \(\mathbf{A}\) 上的一个同构，并把 \(\mathbf{F}\) （相应地 \(\mathbf{E}\) ）看作一个右 \(\mathbf{A}^{0}\) -模。

公式 (24) 与第 3 小节的定义 6 立即表明，对 F（相应地 E）的每个子模 N（相应地 M），有

\[
\mathbf {N ^ {0}} = \stackrel {- 1} {s _ {\Phi}} (\mathbf {N ^ {\prime}}) \quad (\text { resp. } \mathbf {M ^ {0}} = \stackrel {- 1} {d _ {\Phi}} (\mathbf {M ^ {\prime}}))\tag{25}
\]

其中 N′（相应地 M′）是 F 的对偶 F*（相应地 E 的对偶 E*）中与 N（相应地 M）正交的子模（第 II 章，第 4 节，第 2 小节）。

命题 4. --- 设 A 是一个体，\(\Phi\) 是 \(E \times F\) 上的一个双线性型（相应地关于 J 的半双线性型）；\(E/F^{0}\) 是有限维的必须且只需 \(F/E^{0}\) 是有限维的，且这时两者的维数相等。

事实上，设 \(\Phi_{1}\) 为与 \(\Phi\) 相伴的非退化型，定义在 \((\mathrm{E} / \mathrm{F}^{0}) \times (\mathrm{F} / \mathrm{E}^{0})\) 上（第 3 小节）。设 \(\mathrm{E} / \mathrm{F}^{0}\) 是有限维 \(n\) ；由于线性映射 \(d_{\Phi_1}\) （从 \(\mathrm{F} / \mathrm{E}^{0}\) （相应地 \((\mathrm{F} / \mathrm{E}^{0})^{\mathrm{j}})\) ）到 \((\mathrm{E} / \mathrm{F}^{0})^{*}\) 中）是单射，\(\mathrm{F} / \mathrm{E}^{0}\) 是有限维 \(n' \leqslant n\) ；对 \(s_{\Phi_1}\) 作同样考虑，同样看出 \(n \leqslant n'\) 。

推论 1. --- 设 A 是一个体且 \(\Phi\) 非退化。E 的子空间 M 是有限维的，必须且只需 M \(^{0}\) 在 F 中具有有限余维数，且这时有 codim M \(^{0}\) = dim M，以及 M \(^{00}\) = M。

由于 \(F^{0}=\{0\}\) ，前两个断言由把命题 4 应用于 \(\Phi\) 在 M × F 上的限制得出。此外，\(M^{0}\) 是 \(M^{00}\) 的正交补，故 \(M^{00}\) 的维数有限且等于 codim \(M^{0}=dim M\) ；但由于 \(M^{00}\supset M\) ，就有 \(M^{00}=M\) 。

推论 2. --- 在推论 1 的假设下，设 M、N 是 E 的两个子空间；则有 \((M + N)^{0} = M^{0} \cap N^{0}\) ；若此外 M 与 N 都是有限维的，则有 \((M \cap N)^{0} = M^{0} + N^{0}\) 。

第一个断言是显然的。设 M 与 N 都是有限维的，令 \(G = M^{0} + N^{0}\) ；由推论 1 有 \(G^{0} = M^{00} \cap N^{00} = M \cap N\) ；把命题 4 应用于 \(\Phi\) 在 \(M \times G\) 上的限制，便表明（由于 \(M^{0} \subset G\) 且 \(G^{0} \subset M\) ）有 dim \(M/(M \cap N) = \dim G/M^{0} = \text{codim } M^{0} - \text{codim } G\) ，而由 codim \(M^{0} = \dim M\) ，便导出 dim \((M \cap N) = \text{codim } G\) 。但由推论 1，同样有 dim \((M \cap N) = \text{codim } (M \cap N)^{0}\) ，而由于 \(G \subset G^{00} = (M \cap N)^{0}\) ，就有 \(G = (M \cap N)^{0}\) 。

命题 4 使我们可以给出下面的定义：

定义 7. --- 设 A 是一个体（相应地一个配备了反自同构 J 的体），E 是 A 上的一个左向量空间，F 是 A 上的一个右（相应地左）向量空间，Φ 是 E × F 上的一个双线性（相应地关于 J 的半双线性）型。设 E/F⁰ 与 F/E⁰ 在 A 上都是有限维的。Φ 的秩就是向量空间 E/F⁰ 与 F/E⁰ 的公共（有限）维数。

当 \(E/F^{0}\) 与 \(F/E^{0}\) 是无限维时，称 \(\Phi\) 具有无穷秩。

命题 5. --- 在定义 7 的假设与记号下，线性映射 \(s_{\Phi}\) 与 \(d_{\Phi}\) （与 \(\Phi\) 相伴者）有相同的秩，且这个秩等于型 \(\Phi\) 的秩。

事实上，F 到 E* 中的映射 \(d_{\Phi}\) 的核显然是 E°，故其秩等于 F/E° 的维数。同样，\(s_{\Phi}\) 的秩等于 E/F° 的维数。

命题 6. --- 在定义 7 的假设与记号下，再设 E 与 F 有相同的有限维数。则下列条件等价：

\begin{enumerate}
\def\labelenumi{\alph{enumi})}
\item
  \(d_{\Phi}\) 是单射；
\item
  \(d_{\Phi}\) 是满射；
\item
  \(s_{\Phi}\) 是单射；
\end{enumerate}

\(d)\) \(s_{\Phi}\) 是满射；

\begin{enumerate}
\def\labelenumi{\alph{enumi})}
\setcounter{enumi}{4}
\tightlist
\item
  \(\Phi\) 非退化。
\end{enumerate}

事实上，由于 E、F、E* 与 F* 有相同的有限维数，a) 与 b) 等价，c) 与 d) 也等价（第 II 章，第 3 节，第 4 小节）。由于 \(s_{\Phi}\) 与 \(d_{\Phi}\) 有相同的秩（命题 5），a) 与 c) 等价。由于 e) 等价于关系式 \(\mathrm{E}^{0} = \mathrm{F}^{0} = \{0\}\) ，它等价于 a) 与 c) 的合取，从而所述诸条件等价。

推论. --- 在定义 7 的假设与记号下，再设 E 是有限维的且 \(\Phi\) 非退化。则有 dim E = dim F，且对 E 的每个基 \((e_i)\) （ \(1 \leqslant i \leqslant \dim E\) ），存在 F 的一个基 \((f_i)\) ，使得 \(\Phi(e_i, f_k) = \delta_{ik}\) （ \(i, k = 1, \ldots, \dim E\) ）。

事实上，由于 \(\Phi\) 非退化，我们有 \(E^{0} = F^{0} = \{0\}\) ，从而 dim \(E = \dim F\) （命题 4）。由此（命题 6），\(d_{\Phi}\) 是 \(F\) （相应地 \(F^{j}\) ）到 \(E^{*}\) 上的一个同构；因此，若 \((e_{i}^{*})\) 是 \((e_{i})\) 的对偶基，则元素 \(f_{i} = d_{\Phi}^{-1}(e_{i}^{*})\) 构成 \(F\) 的一个基，它鉴于公式 (23)（相应地公式 (24)），满足 \(\Phi(e_{i}, f_{k}) = \delta_{ik}\) 。

显然，在这个推论中，可以把 E 与 F 的作用互换，在证明中把 \(d_{\Phi}\) 换成 \(s_{\Phi}\) 。

注记. --- 设 A 是一个配备了反自同构 J 的环，M 与 N 是右 A-模，\(\Phi\) 是 M × N 上的一个（关于 J 的）左半双线性型（第 2 小节，注记）；于是它满足等式

\[
\Phi (x a, y a ^ {\prime}) = a ^ {j} \Phi (x, y) a ^ {\prime} \quad (a, a ^ {\prime} \in A, x \in M, y \in N).
\]

映射 \(\Phi'\) （由 \(\Phi'(y, x) = \Phi(x, y)^{J'}\) 定义，其中 J′ = J \(^{-1}\) ，从 N × M 到 A 中）是一个关于 J′ 的左半双线性型，且 \(\Phi\) 与 \(\Phi'\) 分别等同于 M \(^{J'}\) × N 与 N \(^{J'}\) × M 上的双线性型。与这两个双线性型相伴的映射 \(s_{\Phi}\) 与 \(s_{\Phi'}\) ，分别称为半双线性型 \(\Phi\) 的左相伴与右相伴映射，记作 \(s_{\Phi}\) 与 \(d_{\Phi}\) 。于是，按定义我们有

\[
\Phi (x, y) = \left\langle y, s _ {\Phi} (x) \right\rangle = \left\langle x, d _ {\Phi} (y) \right\rangle^ {J} \quad (x \in M, y \in N),\tag{26}
\]

且 \(s_{\Phi}\) （相应地 \(d_{\Phi}\) ）是 \(\mathbf{M}^{\mathbf{j}}\) 到 \(\mathbf{N}^{*}\) 中的一个线性映射（相应地，

N \(^{j'}\) 到 M* 中的线性映射）。对这里所考虑的情形，容易叙述并证明定义 7 与命题 4、5、6 的类似结果。

